% Part: propositional-logic
% Chapter: syntax-and-semantics
% Section: introduction

\documentclass[../../../include/open-logic-section]{subfiles}

\begin{document}

\olfileid{pl}{syn}{int}

\olsection{Pambuka}

Logika proposisional ngrembug !!{formula}s kang kawangun saka
!!{propositional variable}s kanthi konektif proposisional
$\lnot$, $\land$, $\lor$, $\lif$, lan $\liff$. Miturut pangerten
lumrah, !!a{propositional variable}~$p$ makili ukara utawa proposisi
kang bener utawa salah. Yen nilai kabeneran saben
!!{propositional variable} ing !!a{formula} wis ditemtokake, nilai
kabeneran saben !!{formula}s kang kawangun saka variabel-variabel iku
nganggo konektif proposisional uga ditemtokake. Kita kandha yen
logika proposisional iku \emph{fungsional tumrap kabeneran}, amarga
semantike ditemtokake dening fungsi-fungsi nilai kabeneran. Mligine,
ing logika proposisional kita ora nggatekake katemtuan liyane ngenani
bener utawa salahe sawijining bab, upamane apa bab iku mesthi bener
lan ora mung bener kanthi kontingen, apa dingerteni yen bab iku bener,
utawa apa bab iku bener saiki, biyen, utawa mbesuk. Kita mung
nggatekake rong nilai kabeneran, yaiku bener
($\True$) lan salah ($\False$), mula ora ngrembug kamungkinan yen
sawijining pratelan ora bener lan ora salah, utawa mung setengah bener.
Kita uga mung nggatekake konektif kang ndadekake nilai kabeneran
!!a{formula} kang dibangun saka konektif kasebut ditemtokake kanthi
jangkep mung dening nilai kabeneran perangan-perangane (lan dudu,
umpamane, dening maknane). Mligine, isih dadi pasulayan apa nilai kabeneran
kondisional ing basa Inggris iku fungsional tumrap kabeneran ing
teges iki. Kondisional material~$\lif$ iku fungsional tumrap kabeneran;
logika liyane ngrembug kondisional kang ora mangkono.

Kanggo ngembangake teori lan metateori logika proposisional kang
fungsional tumrap kabeneran, kita kudu nemtokake luwih dhisik
sintaksis lan semantik ekspresi-ekspresine. Kita bakal njlentrehake salah siji
cara mbangun !!{formula}s saka !!{propositional variable}s nganggo
konektif. Ana uga definisi alternatif. Sistem liyane bisa
milih simbol kang beda, milih himpunan konektif primitif kang beda,
lan nggunakake tandha kurung kanthi cara kang beda (utawa malah ora
nggunakake babar pisan, kaya ing notasi kang diarani notasi Polandia).
Nanging, ing kabeh cara kasebut, aturan pambentukan nemtokake himpunan
!!{formula}s kanthi \emph{induktif}. Yen ditindakake kanthi trep, saben
ekspresi ing pokoke mung bisa kawangun kanthi siji cara miturut aturan
pambentukan. Definisi induktif kang ngasilake ekspresi kang
\emph{mung bisa diwaca kanthi siji cara} ndadekake kita
bisa menehi makna marang ekspresi kasebut kanthi metodhe kang padha---
yaiku definisi induktif.

Menehi makna marang ekspresi iku wewengkone semantik. Konsep utama ing
semantik logika proposisional yaiku panyukupan dening !!a{valuation}.
!!^a{valuation}~$\pAssign{v}$ masangake nilai kabeneran $\True$,
$\False$ marang !!{propositional variable}s. Saben !!{valuation}
nemtokake nilai kabeneran $\pValue{v}(!A)$ kanggo saben
!!{formula}~$!A$. Kanggo !!^a{formula} lan !!a{valuation}~$\pAssign{v}$,
kita kandha yen valuasi kasebut nyukupi formula kasebut yen lan mung yen
$\pValue{v}(!A) = \True$---iki ditulis
$\pSat{v}{!A}$. Relasi iki uga bisa ditemtokake kanthi induksi ing
struktur~$!A$: fungsi kabeneran saka konektif logis digunakake kanggo
nemtokake, umpamane, panyukupan $!A \land !B$ adhedhasar apa $!A$
lan~$!B$ dicukupi utawa ora.

Adhedhasar relasi panyukupan $\pSat{v}{!A}$ kanggo ukara, kita banjur
bisa nemtokake gagasan semantis dhasar yaiku tautologi, akibat
semantis, lan bisa dicukupi. !!^a{formula} iku tautologi,
$\Entails !A$, yen saben !!{valuation} nyukupi formula kasebut, yaiku
$\pValue{v}(!A) = \True$ kanggo saben~$\pAssign{v}$. Formula iku dadi
akibat semantis saka sawijining himpunan !!{formula}s,
$\Gamma \Entails !A$, yen saben !!{valuation} kang nyukupi kabeh
!!{formula}s ing~$\Gamma$ uga nyukupi~$!A$. Lan sawijining himpunan
!!{formula}s bisa dicukupi yen ana !!{valuation} kang bebarengan
nyukupi kabeh !!{formula}s ing himpunan iku. Amarga !!{formula}s
ditemtokake kanthi induktif, lan panyukupan sabanjure uga ditemtokake
kanthi induksi ing struktur !!{formula}s, kita bisa nggunakake induksi
kanggo mbuktekake sipat-sipat semantik kita lan nggayutake gagasan-gagasan
semantis kang wis ditemtokake.


\end{document}
